\section{The Primary Domain Causal DAG and Identification}
\label{sec:meth_domain_dag}

The structural foundation of our causal framework is the \textit{Primary Domain Causal DAG} ($\mathcal{G}_{\text{domain}}$), constructed in consultation with higher education pedagogical literature and institutional regulations at IPPortalegre \parencite{martins2021predicting}.

\subsection{Node Taxonomy and Operational Definitions}

The 35 administrative features are mapped into a 7-node canonical DAG:
\begin{enumerate}
    \item \textbf{Socioeconomic Profile ($Z_1$):} Student enrollment age, parental educational attainment, and housing displacement status.
    \item \textbf{Prior Academics ($Z_2$):} Secondary school exit grade and national entrance examination performance (\texttt{Admission grade}).
    \item \textbf{Financial Solvency ($Z_3$):} Baseline tuition fee arrears status upon enrollment (\texttt{Debtor}).
    \item \textbf{Primary Policy Treatment ($T$):} Binary indicator for state-subsidized operational scholarship award ($T=1$: Awarded, $T=0$: Control).
    \item \textbf{Semester 1 Credits ($M_1$):} Total ECTS course units approved in first semester (\texttt{Curricular units 1st sem approved}).
    \item \textbf{Semester 1 GPA ($M_2$):} Normalized first-semester grade point average ($[0, 20]$ scale).
    \item \textbf{Terminal Outcome ($Y$):} Definitive educational outcome ($Y=0$: Dropout, $Y=1$: Graduate).
\end{enumerate}

\subsection{Graphical Topology and Backdoor Admissibility}

Figure \ref{fig:domain_dag_tikz} displays the causal topology $\mathcal{G}_{\text{domain}}$.

\begin{figure}[htbp]
    \centering
    \resizebox{0.82\textwidth}{!}{%
    \begin{tikzpicture}[
        node distance=1.6cm and 2.2cm,
        every node/.style={draw, font=\small\sffamily, thick},
        confounder/.style={rectangle, rounded corners=3pt, fill=blue!10, text width=2.6cm, align=center},
        treatment/.style={rectangle, rounded corners=3pt, fill=red!15, draw=red!70!black, line width=1.5pt, text width=2.4cm, align=center},
        mediator/.style={rectangle, rounded corners=3pt, fill=orange!15, draw=orange!80!black, line width=1.2pt, text width=2.6cm, align=center},
        outcome/.style={rectangle, rounded corners=3pt, fill=green!20, draw=green!60!black, line width=1.5pt, text width=2.4cm, align=center},
        edge/.style={->, >=Latex, thick, draw=black!80},
        causal_edge/.style={->, >=Latex, very thick, draw=red!80!black},
        med_edge/.style={->, >=Latex, thick, draw=orange!80!black}
    ]
        % Top row: Confounders
        \node[confounder] (Z1) at (0, 2.6) {Socioeconomic ($Z_1$)};
        \node[confounder] (Z2) at (4.2, 2.6) {Prior Academics ($Z_2$)};
        \node[confounder] (Z3) at (8.4, 2.6) {Financial Solvency ($Z_3$)};
        
        % Middle row: Treatment & Mediators
        \node[treatment] (T) at (0, 0) {\textbf{Scholarship ($T$)}};
        \node[mediator] (M1) at (4.2, 0.7) {Sem 1 Credits ($M_1$)};
        \node[mediator] (M2) at (4.2, -0.7) {Sem 1 GPA ($M_2$)};
        
        % Bottom right: Outcome
        \node[outcome] (Y) at (8.4, 0) {\textbf{Graduation ($Y$)}};
        
        % Confounder Edges
        \draw[edge] (Z1) -- (T);
        \draw[edge] (Z1) -- (M1);
        \draw[edge] (Z1) to[out=-40, in=140] (Y);
        
        \draw[edge] (Z2) -- (T);
        \draw[edge] (Z2) -- (M1);
        \draw[edge] (Z2) -- (M2);
        \draw[edge] (Z2) -- (Y);
        
        \draw[edge] (Z3) -- (T);
        \draw[edge] (Z3) -- (Y);
        
        % Causal Paths
        \draw[causal_edge] (T) to[out=-30, in=-150] node[below, font=\scriptsize, draw=none, fill=none] {Direct Effect ($NDE$)} (Y);
        \draw[med_edge] (T) -- (M1);
        \draw[med_edge] (T) -- (M2);
        \draw[med_edge] (M1) -- (Y);
        \draw[med_edge] (M2) -- (Y);
        \draw[edge] (M1) -- (M2);
    \end{tikzpicture}%
    }
    \caption{The Primary Domain Causal DAG ($\mathcal{G}_{\text{domain}}$). Red node indicates treatment; orange nodes denote academic mediators; green node is terminal graduation.}
    \label{fig:domain_dag_tikz}
\end{figure}

\begin{theorem}[Backdoor Identification of Scholarship Effect]
Under $\mathcal{G}_{\text{domain}}$, the conditioning set $\mathbf{Z} = \{Z_1, Z_2, Z_3\}$ satisfies the Backdoor Criterion relative to $(T, Y)$. Consequently, $P(Y \mid do(T=t))$ is identified non-parametrically by:
\begin{equation}
    P(Y=1 \mid do(T=t)) = \sum_{z_1, z_2, z_3} P(Y=1 \mid T=t, z_1, z_2, z_3) P(z_1, z_2, z_3).
    \label{eq:backdoor_proof_domain}
\end{equation}
\end{theorem}
\begin{proof}
First, no node in $\mathbf{Z}$ is a descendant of $T$; all elements of $\mathbf{Z}$ represent pre-enrollment baseline traits determined strictly prior to scholarship allocation. Second, in the manipulated sub-graph $\mathcal{G}_{\underline{T}}$, all outgoing arrows from $T$ ($T \to M_1, T \to M_2, T \to Y$) are removed. The remaining backdoor paths connecting $T$ and $Y$ are: $T \leftarrow Z_1 \to Y$, $T \leftarrow Z_2 \to Y$, and $T \leftarrow Z_3 \to Y$. Since conditioning on $\mathbf{Z} = \{Z_1, Z_2, Z_3\}$ blocks all three forks, d-separation holds: $(T \perp_{\mathcal{G}_{\underline{T}}} Y \mid \mathbf{Z})$.
\end{proof}


\subsection{Formal Validity Constraints for Causal DAGs}
To audit whether candidate graph topologies generated by automated discovery algorithms or proposed by human domain analysts are physically and mathematically valid, we enforce a four-pillar constraint verification protocol:
\begin{enumerate}
    \item \textbf{Acyclicity Constraint:} The candidate graph $\mathcal{G} = (\mathbf{V}, \mathbf{E})$ must contain strictly zero directed cycles: $\operatorname{tr}(e^{W \odot W}) - d = 0$. In algorithmic implementations, this is audited via Tarjan's cycle detection or topological sort.
    \item \textbf{Temporal Tier Consistency:} Features are partitioned into four non-invertible administrative stages: Tier 1 (Pre-enrollment $Z_1, Z_2, Z_3$) $\prec$ Tier 2 (Admission Aid $T$) $\prec$ Tier 3 (First Semester $M_1, M_2$) $\prec$ Tier 4 (Graduation Outcome $Y$). Any candidate edge directing causality backward in time belongs to the forbidden set $\mathcal{E}_{\text{forbidden}} = \{(u, v) \mid \text{tier}(u) > \text{tier}(v)\}$ and is rejected.
    \item \textbf{Domain Blacklist and Whitelist Adherence:} Must strictly include substantive institutional pathways ($T \to Y$, $T \to M_1$, $M_1 \to Y$) while strictly prohibiting absurd feedback arrows ($Y \to T$, $M_1 \to Z_1$).
    \item \textbf{Backdoor $d$-Separation Admissibility:} On the manipulated subgraph $\mathcal{G}_{\underline{T}}$ where all outgoing edges from $T$ are pruned, the candidate conditioning set $\mathbf{Z}$ must d-separate $T$ from $Y$: $(T \perp_{\mathcal{G}_{\underline{T}}} Y \mid \mathbf{Z})$.
\end{enumerate}
